\documentclass{beamer}

\usepackage{xcolor}
\usepackage{tikz}

% Theme choice
\usetheme{Boadilla}
% 120,180,250
\definecolor{MyBlue}{RGB}{0,40,255}
\usecolortheme[named=MyBlue]{structure}

% Packages
\usepackage[utf8]{inputenc}
\usepackage{graphicx}

% Title page info
\title{Differenzierung TODO}
\author{Noah Albers, Felix Graff}
\institute{Uni Münster}
\date{\today}


\usepackage{comment}

\newcommand{\teacher}[1]{\textcolor{red}{#1}}
\newcommand{\tteacher}[1]{#1}
%\newcommand{\ppause}{\pause}

%\newcommand{\tteacher}[1]{}
%\newcommand{\teacher}[1]{}
\newcommand{\ppause}{}


\newcommand{\lowlight}[1]{\textcolor{gray}{#1}}



\begin{document}
	
	% Title slide
	\frame{\titlepage}
	
	
	\begin{frame}{Übersicht}
		\tableofcontents
	\end{frame}

	\begin{frame}{Was ist Scratch?}{Die Idee dahinter}
		\begin{itemize}
			\item Blockbasiertes Interface zum erstellen kleiner Stories
			\item 2007 bei MIT Media Lab entwickelt
			\item 2019 Übergang in die "Scratch Foundation" (non profit)
			\item Verschiedene Möglichkeiten:
			\begin{itemize}
				\item Vorgemalte Charaktere, Landschaften, Sounds, Grafikten etc.
				\item Eigene Elemente bzw. Bilder hochladen
				\item In Scratch selber erstellen/bearbeiten
			\end{itemize}
			\item Entwickelt für Kinder
			\begin{itemize}
				\item Programmieren spielerisch erlernen 
				\item Mathematischer visualisieren
				\item Ideen freien Lauf lassen
			\end{itemize}
		\end{itemize}
		
	\end{frame}


	\begin{frame}{Unterrichtsszenario}{Rahmenbedingungen}
		\begin{itemize}
			\item Fach: Informatik
			\item Klasse: Sekundarstufe I (5-6)
			\item Thema: Sequenzielles Abarbeiten (Algorithmik) \\
			(Durch faded worked examples)
			\item Ziele:
			\begin{itemize}
				\item SuS wissen, dass Programme Schrittweise ablaufen 
				\item SuS wissen, wie einfache Schleifen funktionieren
				\item SuS können blockbasierte Programme lesen/modifizieren/erstellen
			\end{itemize}
		\end{itemize}
		
	\end{frame}


	\begin{frame}{Scratch-Beispiel}
		
	\end{frame}
	
	
	
\end{document}